\scanpage{132}
A un tel $L$, on associe la structure complexe de \deleted{$V$ déduite \inserted{\deleted{\unsure{conjuguée}}} de celle \inserted{\deleted{\unsure{déduite}}} de $L$ par transport de} \deleted{\unsure{conjuguée} de celle déduite de} $V$ qui est conjuguée de celle déduite de la str. complexe de $L$ par transp de structure via (7), i.e. on pose
\[
 ux=\varphi(-i\varphi^{-1}(x))=-\varphi(i\varphi(x)) \tag*{(8)}
\]
$u$ étant la multiplication par $i$ pour la structure complexe envisagée sur $V$.

% slanted
\margin{cette convention est imposée par la situation de ``modules'' des courbes elliptiques et du fait qu'on veut que les paramètres qu'on \unsure{(va)} y introduit soient \uline{holomorphes}, et non antiholomorphes \ldots}

Inversement, si $V$ est muni d'une structure complexe, \deleted{l'applic} faisant de $V$ un espace vectoriel complexe $V_c$, alors l'application \inserted{identique} $\mathbb R$-linéaire $V\longrightarrow V_c$ se prolonge, par la propriété universelle de $V_{\mathbb C}$, en une application $\mathbb C$-linéaire
\[
 \psi\colon V_{\mathbb C}\longrightarrow V_c \tag*{(9)}
\]
de sorte que
\[
 \left\{\begin{aligned}\psi(ix)&=ux&&\forall x\in V=V_{\mathbb R},\\ \psi(x)&=x\end{aligned}\right. \tag*{(10)}
\]
donc
\[
 \psi(x+iy)=x+uy \tag*{(11)}
\]
\deleted{Considérons} \inserted{\deleted{\unsure{Posons|Prenons}}} Considérons
\[
 \begin{aligned}
 L^{\unsure{\prime|1}}=\mathrm{Ker}\,\psi
 &=\{x+iy\mid x,y\in V_{\mathbb R},\ x+uy=0\quad\text{i.e. }y=\deleted{\unsure{-}}ux\}\\
 &=\{\deleted{\unsure{uy+iy}}\inserted{x+iu(x)}\mid x\in V_{\mathbb R}\}
 \end{aligned}\tag*{(12)}
\]
et considérons l'application correspondante (6)
\[
 \varphi\colon\underset{\substack{\Vert\\x+iu(x)}}{z}\longmapsto\tfrac12(z+\bar z)=\deleted{\unsure{x+uy}}\inserted{\Re z=x}\colon L\longrightarrow V
\]
On voit alors que cette application est bijective, et \uline{antilinéaire}, \inserted{donc \deleted{la} correspond à la structure complexe} \deleted{\unsure{Prenant une|Pour un} \unsure{pour $V$ de dim. finie} sur $\mathbb R$, $L=L'=\overline{\mathrm{Ker}\,\psi}$ \unsure{$\subset V_{\mathbb C}$}, \unsure{$\mathrm{dim}_{\mathbb R}=2|\mathrm{rang}_{\mathbb R}=2$},}
% slanted
\margin{mise sur $V$. alors $L\longrightarrow V_c$ est \uline{\unsure{linéaire|antilinéaire}}}

\[
 \tag*{(13)}
 \left[\begin{minipage}{0.84\linewidth}
 On peut dire aussi, que la donnée d'une structure complexe sur $V$ équivaut à

 $1^\circ$) La donnée d'un \deleted{\unsure{opé}} sous-groupe compact \deleted{maximal} $U$ \deleted{de $\mathrm{Gl}=$} $\inserted{G=}\mathrm{Aut}_{\mathbb R}(W)$, \deleted{i.e. d'un}
 \end{minipage}\right.
\]
% archive sheet number: 131
